\documentclass[11pt]{article}

% ---------------------------------------------------------------------------
%  Companion notes for the interactive demos at matteovissani.github.io
%  Mathematical model + numerical scheme + simulated data for each lab.
% ---------------------------------------------------------------------------

\usepackage[a4paper,margin=1in]{geometry}
\usepackage{amsmath,amssymb,mathtools}
\usepackage{xcolor}
\usepackage{titlesec}
\usepackage{enumitem}
\usepackage[colorlinks=true,allcolors=mygreen]{hyperref}
\usepackage{parskip}
\usepackage{booktabs}

% --- synthwave-ish accent palette (kept subtle on a white page) -----------
\definecolor{neoncyan}{HTML}{0E8FA8}
\definecolor{neonmag}{HTML}{C81E6E}
\definecolor{violet}{HTML}{5B4BC4}
\definecolor{mygreen}{HTML}{1E8C5A}
\definecolor{rule}{HTML}{C81E6E}

\titleformat{\section}{\Large\bfseries\color{violet}}{\thesection.}{0.6em}{}
\titleformat{\subsection}{\large\bfseries\color{neoncyan}}{\thesubsection}{0.6em}{}
\titlespacing{\section}{0pt}{1.6em}{0.6em}

\newcommand{\param}[1]{\textcolor{neonmag}{\texttt{#1}}}
\newcommand{\code}[1]{\texttt{#1}}
\newcommand{\E}{\mathrm{e}}
\newcommand{\ii}{\mathrm{i}}
\newcommand{\dd}{\mathrm{d}}

\setlist[itemize]{leftmargin=1.4em,itemsep=2pt,topsep=3pt}

\begin{document}

% ======================= TITLE =======================
\begin{titlepage}
\centering
\vspace*{3cm}
{\color{rule}\rule{\textwidth}{2pt}}\\[1.2em]
{\Huge\bfseries Interactive Demos}\\[0.4em]
{\LARGE Mathematical \& Simulation Notes}\\[1em]
{\color{rule}\rule{\textwidth}{2pt}}\\[2em]
{\large A companion to the interactive labs at}\\[0.3em]
{\large\href{https://matteovissani.github.io}{matteovissani.github.io}}\\[3em]
{\large\bfseries Matteo Vissani}\\[0.3em]
{\normalsize Neuroengineer --- Massachusetts General Hospital / Harvard Medical School}\\[2em]
\vfill
{\small Every figure on the site is computed live in the browser. These notes give, for
each lab, the exact model, the numerical scheme, the synthetic data it runs on, and the
quantities it measures --- so the demos can be read, checked, and reused as teaching material.}\\[1em]
{\small\itshape Generated \today.}
\end{titlepage}

\tableofcontents
\newpage

% =====================================================================
\section*{How to read these notes}
\addcontentsline{toc}{section}{How to read these notes}

Each lab is a small, self-contained simulation that runs in real time on an HTML5 canvas
(no server, no precomputed data). For every lab below you will find:

\begin{itemize}
  \item \textbf{Model} --- the governing equations.
  \item \textbf{Numerics} --- the integration / estimation scheme and step sizes.
  \item \textbf{Synthetic data} --- how the signals or spikes are generated.
  \item \textbf{Measured} --- the quantities reported back to the user.
  \item \textbf{Controls} --- the adjustable parameters and their ranges.
\end{itemize}

Parameters that appear as on-screen sliders are written like \param{this}. The labs are
grouped exactly as on the website: \emph{neuronal dynamics}, \emph{neural coding \&
information}, \emph{signals \& time--frequency}, \emph{neurotechnology}, and the
\emph{spike--phase coupling toolkit}.

% #####################################################################
\section{Neuronal dynamics}

% ---------------------------------------------------------------------
\subsection{The (adaptive) integrate-and-fire neuron}

\textbf{Model.} A leaky integrate-and-fire neuron with an adaptation current $w$:
\begin{align}
\tau_m \frac{\dd V}{\dd t} &= -(V - E_L) + R I - w + \text{noise}, &
\tau_w \frac{\dd w}{\dd t} &= -w,
\end{align}
with the fire-and-reset rule: when $V \ge V_{\text{th}}$ emit a spike, set $V \leftarrow V_r$
and $w \leftarrow w + b$. Fixed constants are $E_L=-70$, $V_{\text{th}}=-50$, $V_r=-65$~mV.
Setting $b=0$ recovers the pure LIF neuron; $b>0$ produces spike-frequency adaptation.

\textbf{Numerics.} Euler--Maruyama with $\Delta t = 0.5$~ms. The stochastic term is added as
$\sigma\sqrt{\Delta t/\tau_m}\,\xi$ with $\xi\sim\mathcal N(0,1)$ (Box--Muller). The membrane
panel scrolls an $800$~ms ring buffer; the integrator advances $8$ substeps per animation frame.

\textbf{Phase plane.} The $(V,w)$ plane shows the two nullclines
\[
\text{$V$-nullcline:}\quad w = R I - (V - E_L), \qquad
\text{$w$-nullcline:}\quad w = 0,
\]
with the live trajectory riding up toward threshold, resetting leftward, and being kicked up
by $b$ at each spike.

\textbf{$f$--$I$ curve.} For the current shape parameters the gain curve is measured by a
direct sweep of the drive over $[0,45]$~mV: each point integrates a $400$~ms burn-in followed
by a $1400$~ms count window ($\Delta t = 0.5$~ms) and reports the firing rate
$f = 1000\,n_{\text{spikes}}/T$ (Hz). With $w=0$ the steady state is $V=E_L+RI$, so the
\emph{rheobase} (minimum drive to fire) is $V_{\text{th}}-E_L = 20$~mV.

\textbf{Controls.} \param{drive $RI$} $\in[0,45]$~mV, \param{noise $\sigma$} $\in[0,8]$~mV,
\param{adaptation $b$} $\in[0,10]$~mV, \param{$\tau_w$} $\in[40,300]$~ms,
\param{$\tau_m$} $\in[5,40]$~ms.

% ---------------------------------------------------------------------
\subsection{Wilson--Cowan dynamics \& bifurcation}

\textbf{Model.} Two coupled populations, excitatory $E$ and inhibitory $I$:
\begin{align}
\tau_E \dot E &= -E + S\!\left(w_{EE} E - w_{EI} I + P\right), &
\tau_I \dot I &= -I + S\!\left(w_{IE} E - w_{II} I + Q\right),
\end{align}
with the logistic gain $S(x) = S_{\max}\big/\big(1 + \E^{-a(x-\theta)}\big)$,
$S_{\max}=1$, $\theta=3$. Fixed: $\tau_E=1$, $w_{IE}=10$, $w_{II}=2$, $Q=0$.

\textbf{Numerics.} Forward Euler, $\Delta t = 0.04$. The phase-plane flow field samples
$f(E,I)=(\dot E,\dot I)$ on an $11\times11$ grid (arrows), and tracer particles are advected
along $f$ to visualize the basins of attraction.

\textbf{Nullclines and fixed points.} Inverting the sigmoid,
$S^{-1}(y)=\theta + \frac1a\ln\!\frac{y}{S_{\max}-y}$, gives the closed-form nullclines
\[
\text{$E$-null:}\;\; I = \frac{w_{EE}E + P - S^{-1}(E)}{w_{EI}},\qquad
\text{$I$-null:}\;\; E = \frac{S^{-1}(I) + w_{II}I - Q}{w_{IE}}.
\]
Fixed points are the roots of $g(E) = -E + S\!\big(w_{EE}E - w_{EI} I^\ast(E) + P\big)$,
where $I^\ast(E)$ is obtained by fixed-point iteration of the $I$ equation; a sign-change scan
locates the $1$ or $3$ intersections (single state, or bistable cusp / saddle-node regime).

\textbf{Bifurcation diagram.} The external drive is swept $P\in[0,12]$ ($80$ values). For each
$P$ the system settles ($1600$ steps), then the running min/max of $E$ over the next $1500$
steps gives the lower/upper branch; their separation is the oscillation amplitude, which grows
continuously from zero at a \emph{supercritical Hopf bifurcation}.

\textbf{Controls.} \param{$P$} $\in[0,12]$, \param{$w_{EE}$} $\in[4,16]$,
\param{$w_{EI}$} $\in[6,22]$, \param{$\tau_I$} $\in[1,4]$, \param{gain $a$} $\in[0.8,2.5]$.

% #####################################################################
\section{Neural coding \& information}

% ---------------------------------------------------------------------
\subsection{Spike--phase coupling (teaching demo)}

\textbf{Model.} A cortical rhythm $\cos$ wave of frequency $f$ drives a subthalamic unit whose
instantaneous firing rate is phase-modulated by a von Mises-type kernel,
\[
\lambda(\phi) = r_0\, \exp\!\big(\kappa[\cos(\phi-\phi_{\text{pref}})-1]\big),
\]
bounded in $(0,1]\cdot r_0$, where $\phi = 2\pi f t$ is the oscillation phase, $\kappa$ the
concentration (coupling strength) and $\phi_{\text{pref}} = 2\pi f \cdot(\text{offset})$.

\textbf{Numerics.} Inhomogeneous Poisson sampling at $\Delta t = 1$~ms: a spike is emitted when
$\mathcal U(0,1) < \lambda(\phi)\,\Delta t$. The last $180$ spike phases are kept in a ring buffer.

\textbf{Measured.} The phase-locking value and preferred phase from the unit resultant vector,
\[
\text{PLV} = \Big|\tfrac1n\textstyle\sum_{k=1}^{n}\E^{\ii\phi_k}\Big|, \qquad
\phi_{\text{pref}} = \arg\!\Big(\textstyle\sum_k \E^{\ii\phi_k}\Big),
\]
with an $18$-bin polar histogram and the mean vector drawn at length PLV.

\textbf{Controls.} \param{frequency} $\in[2,40]$~Hz ($\delta$ to $\gamma$),
\param{$\kappa$} $\in[0,5]$, \param{phase offset} $\in[0,40]$~ms,
\param{firing rate} $\in[5,80]$~Hz.

% ---------------------------------------------------------------------
\subsection{Population coding \& neural decoding}

\textbf{Model.} A bank of $N$ Gaussian tuning curves over a stimulus $s\in[0,180]^\circ$,
\[
f_i(s) = 0.5 + r_{\max}\,\exp\!\Big(-\frac{(s-c_i)^2}{2\sigma^2}\Big), \qquad
c_i = \frac{i}{N-1}\,180^\circ,
\]
producing independent Poisson spike counts $r_i\sim\text{Poisson}(f_i(s))$ each trial.

\textbf{Decoder.} Maximum likelihood over a $1^\circ$ grid; dropping $s$-independent terms,
\[
\hat s = \arg\max_{g}\ \sum_{i=1}^{N}\big[\,r_i \ln f_i(g) - f_i(g)\,\big].
\]

\textbf{Optimality bound.} For Poisson firing with Gaussian tuning the Fisher information is
\[
J(s) = \sum_{i} \frac{f_i(s)\,(c_i-s)^2}{\sigma^4},
\]
and the Cram\'er--Rao bound on any unbiased estimator's standard deviation is
$\mathrm{CRB} = 1/\sqrt{J}$. The lab overlays this Gaussian on the empirical error histogram
(last $200$ trials) and reports the ratio $\mathrm{RMSE}/\mathrm{CRB}$ (near $1$ $\Rightarrow$
near-optimal). Poisson counts use direct multiplication for small $\lambda$ and a Gaussian
approximation for $\lambda>30$.

\textbf{Controls.} \param{$s$} $\in[10,170]^\circ$, \param{$\sigma$} $\in[6,45]^\circ$,
\param{$N$} $\in[6,40]$, \param{$r_{\max}$} $\in[5,50]$.

% ---------------------------------------------------------------------
\subsection{Information theory: channel \& source coding}

\textbf{Part 1 --- binary symmetric channel.} A bit $X$ with $P(X{=}1)=p$ is flipped with
probability $f$. With the binary entropy $H_2(q) = -q\log_2 q - (1-q)\log_2(1-q)$,
\[
H(X)=H_2(p),\quad H(Y\,|\,X)=H_2(f),\quad p_Y = p(1-f)+(1-p)f,
\]
\[
I(X;Y) = H(Y) - H(Y\,|\,X), \qquad C = 1 - H_2(f),
\]
and the joint entropy decomposes as $H(X,Y) = H(X\,|\,Y) + I(X;Y) + H(Y\,|\,X)$, shown as a
stacked bar. The mutual information is maximal at $p=0.5$, where it equals the capacity $C$.

\textbf{Part 2 --- source coding.} For a symbol distribution $\{p_i\}$ the entropy
$H = -\sum_i p_i \log_2 p_i$ is the compression limit. An exact \textbf{Huffman} code is built
by repeatedly merging the two least-probable nodes; its expected length $L=\sum_i p_i \ell_i$
obeys Shannon's bound $H \le L < H+1$. The lab reports the efficiency $H/L$, which reaches
$100\%$ when the probabilities are (negative) powers of two.

\textbf{Controls.} \param{$p$} $\in[0,1]$, \param{$f$} $\in[0,0.5]$, and five symbol
\param{weights} $\in[0,8]$.

% #####################################################################
\section{Signals \& time--frequency}

% ---------------------------------------------------------------------
\subsection{Fourier epicycles}

\textbf{Model.} A closed planar curve is sampled to $N=220$ complex points
$z_n = x_n + \ii y_n$ (mean-centered) and transformed by the DFT,
\[
c_k = \frac1N \sum_{n=0}^{N-1} z_n\, \E^{-\ii 2\pi kn/N}, \qquad
z(t) = \sum_{k} c_k\, \E^{\ii 2\pi f_k t},
\]
where each coefficient is one rotating vector (epicycle) of radius $|c_k|$, angular frequency
$f_k$ (mapped to $[-N/2, N/2]$) and phase $\arg c_k$. Coefficients are sorted by amplitude and
the slider keeps the largest $K$; chained tip-to-tail they retrace the shape.

\textbf{Numerics.} The user-drawn or preset path is arc-length resampled to $N$ equal steps
before the DFT, so the parametrization speed is uniform. Sharp corners require many
high-frequency terms (the square-wave Gibbs phenomenon).

\textbf{Controls.} \param{number of terms} $\in[1,220]$, \param{speed} $\in[0.2,4]\times$;
or draw a custom shape on the canvas.

% ---------------------------------------------------------------------
\subsection{Digital filter designer (windowed-sinc FIR)}

\textbf{Model.} A band-pass FIR filter from the windowed-sinc method at $F_s = 250$~Hz. With
normalized edges $\omega_{L,H} = 2\pi f_{L,H}/F_s$ and $M$ taps centred at $\text{mid}=(M{-}1)/2$,
\[
h_{\text{ideal}}[n] =
\begin{cases}
(\omega_H - \omega_L)/\pi, & n = 0,\\[4pt]
\dfrac{\sin(\omega_H n) - \sin(\omega_L n)}{\pi n}, & n \ne 0,
\end{cases}
\qquad h[k] = h_{\text{ideal}}[k-\text{mid}]\cdot w[k],
\]
with a Hamming window $w[k] = 0.54 - 0.46\cos\!\big(2\pi k/(M-1)\big)$.

\textbf{Response \& application.} The magnitude response is evaluated directly,
\[
|H(f)| = \Big|\textstyle\sum_k h[k]\,\E^{-\ii 2\pi f k / F_s}\Big|,
\]
and the filter is applied by centred convolution. The test signal is
$\sin(2\pi 8t) + \sin(2\pi 40 t) + 0.6\sin(2\pi 60 t) + 0.5\,\eta(t)$,
where $\eta$ is $1/f$-like noise from an AR(1) process ($a=0.92$). Input and output power spectra
are compared by direct DFT. More taps $\Rightarrow$ sharper passband but a longer impulse
response (group delay) --- the FIR length vs.\ selectivity trade-off.

\textbf{Controls.} \param{low edge} $\in[1,120]$~Hz, \param{high edge} $\in[3,124]$~Hz,
\param{taps} $\in[11,201]$ (forced odd); Nyquist $=125$~Hz.

% ---------------------------------------------------------------------
\subsection{Wavelet time--frequency decomposition}

\textbf{Model.} A complex Morlet continuous wavelet transform. The wavelet at analysis
frequency $f$ is a complex exponential under a Gaussian envelope,
\[
\psi_f(\tau) = \E^{\,\ii 2\pi f \tau}\,\E^{-\tau^2/(2\sigma^2)}, \qquad
\text{CWT}(f,t) = \sum_\tau x(t+\tau)\,\psi_f^{*}(\tau).
\]
The width $\sigma$ (seconds) can be set three equivalent ways (Cohen, 2019), with
$K = 2\sqrt{2\ln 2}$ converting a Gaussian std to its FWHM:
\[
\underbrace{\sigma = \frac{n_{\text{cyc}}}{2\pi f}}_{\text{constant-}Q},\qquad
\underbrace{\sigma = \frac{\text{FWHM}_t}{K}}_{\text{fixed time res.}},\qquad
\underbrace{\sigma = \frac{1}{2\pi(\text{FWHM}_f/K)}}_{\text{fixed freq.\ res.}}
\]
The two resolutions always obey the uncertainty relation
$\text{FWHM}_t \cdot \text{FWHM}_f \approx K^2/(2\pi) \approx 0.88$.

\textbf{Numerics.} $F_s = 200$~Hz, $3$~s. The scalogram uses $46$ log-spaced frequencies
($2$--$50$~Hz); power is $\log_{10}(|\text{CWT}|^2/(\sigma_{\text{samp}}{+}1))$ then
normalized for display. The test signal is a $\theta$ rhythm $+$ a Gaussian-windowed $\beta$
burst (centre $1.5$~s, $\sigma_{\text{win}}=0.28$~s) $+$ $1/f$ noise.

\textbf{Controls.} wavelet-width mode (\param{$n_{\text{cyc}}$} $\in[2,15]$ /
\param{FWHM time} $\in[40,800]$~ms / \param{FWHM freq} $\in[1,15]$~Hz),
$\theta$ and $\beta$ \param{frequency}/\param{amplitude}, \param{$1/f$ noise}, burst on/off.

% ---------------------------------------------------------------------
\subsection{Kalman-filter BCI decoder}

\textbf{Model.} A constant-velocity (near-constant-velocity) Kalman filter runs independently
per axis with state $\mathbf x = [\,\text{pos},\ \text{vel}\,]^\top$. The intended cursor follows
a spring--damper toward the active center-out target, $\ddot p = 6.5(p_{\text{tgt}}-p) - 3.2\,\dot p$,
and the ``neural read-out'' is a noisy position measurement $z = p_{\text{true}} + \sigma\,\xi$.

\textbf{Filter.} With $\Delta t = 0.045$~s, measurement variance $R=\sigma^2$ and process
(acceleration) intensity $q$, the predict step propagates the covariance
$\mathbf P=\big[\begin{smallmatrix}P_{00}&P_{01}\\P_{01}&P_{11}\end{smallmatrix}\big]$ with the
continuous white-noise-acceleration model:
\begin{align*}
P_{00} &\mathrel{+}= 2\Delta t\,P_{01} + \Delta t^2 P_{11} + \tfrac13 q\,\Delta t^3, &
P_{01} &\mathrel{+}= \Delta t\,P_{11} + \tfrac12 q\,\Delta t^2, &
P_{11} &\mathrel{+}= q\,\Delta t,
\end{align*}
and the (scalar-measurement) update uses gain $K_0 = P_{00}/S$, $K_1 = P_{01}/S$ with
$S = P_{00}+R$ and innovation $z - \hat p$. The lab compares the RMSE of the raw read-out and
the Kalman estimate over the last $200$ steps. Too small $q$ lags fast reaches (over-smoothing);
too large $q$ chases noise (under-smoothing) --- the BCI tuning problem.

\textbf{Controls.} \param{measurement noise $\sigma$} $\in[0.03,0.45]$,
\param{process noise $q$} $\in[0.001,0.1]$.

% #####################################################################
\section{Neurotechnology}

% ---------------------------------------------------------------------
\subsection{DBS waveform \& charge-safety designer}

\textbf{Model.} A deep-brain-stimulation pulse is parametrized by amplitude $I$ (mA), pulse
width $\text{PW}$ ($\mu$s), frequency, electrode area $A$, and waveform (monophasic, biphasic
symmetric, biphasic asymmetric with a $4\times$ longer / $\tfrac14$ amplitude recharge phase
and a $60~\mu$s interphase gap). The charge engineering is
\[
Q = \frac{I\cdot\text{PW}}{1000}\ [\mu\text{C/phase}],\qquad
D = \frac{Q}{A/100}\ [\mu\text{C/cm}^2],\qquad
\bar I = Q\cdot f\ [\mu\text{A}].
\]

\textbf{Safety.} The Shannon criterion for tissue damage uses
$k = \log_{10} D + \log_{10} Q$; the plot draws the iso-$k$ lines $\log_{10} D = k - \log_{10} Q$
at $k=1.5$ (safe boundary) and $k=2.0$ (damage). The verdict is SAFE ($k<1.5$),
CAUTION ($1.5\le k<2.0$) or UNSAFE ($k\ge2.0$). Monophasic pulses leave an uncompensated net DC
charge $Q$ per pulse (electrode corrosion / tissue damage); charge-balanced biphasic pulses net zero.

\textbf{Controls.} \param{$I$} $\in[0.1,6]$~mA, \param{PW} $\in[20,450]~\mu$s,
\param{frequency} $\in[2,250]$~Hz, \param{area} $\in[1,12]$~mm$^2$, waveform type.

% ---------------------------------------------------------------------
\subsection{EEG source localization}

\textbf{Forward model.} An illustrative 2-D head: each cortical current dipole
$s=(x,y,\text{ori})$ contributes a potential to a sensor at $(q_x,q_y)$,
\[
V(q; s) = \frac{\cos(\text{ori})\,\Delta x + \sin(\text{ori})\,\Delta y}{\Delta x^2 + \Delta y^2 + \varepsilon},
\quad (\Delta x,\Delta y) = (q_x - x,\ q_y - y),\ \ \varepsilon = 0.0035,
\]
and the measurement at electrode $i$ of the $19$-channel $10$--$20$ montage is
$y_i = \sum_{\text{sources}} V(\text{elec}_i; s) + \text{noise}\cdot\xi_i$ (volume conduction =
spatial mixing). The scalp topomap is a Shepard inverse-distance interpolation, weights
$w = 1/(d^2 + 0.02)$.

\textbf{Inverse.} A greedy / matching-pursuit $K$-dipole fit: on each pass the best grid location
is found by solving the $2\times2$ normal equations for the dipole moment $(p_0,p_1)$ that
minimizes the residual
\[
\sum_i\Big(\text{resid}_i - \tfrac{\Delta x\,p_0 + \Delta y\,p_1}{r^2}\Big)^2;
\]
that dipole is then peeled off the residual and the search repeats. The reported error is the
mean distance (scaled to mm) from each true source to its nearest fitted dipole. Under-fitting,
deep sources, or high noise destabilize the estimate --- the EEG inverse problem is ill-posed.

\textbf{Controls.} drag dipoles; \param{orientation} $\in[0,360]^\circ$,
\param{sensor noise} $\in[0,0.5]$, number of \param{fit dipoles} $\in\{1,2,3,4\}$.

% #####################################################################
\section{Spike--phase coupling toolkit}

This is a browser port of the analysis pipeline from
\href{https://github.com/Brain-Modulation-Lab/code_SPC_ECoG_STN_Speech}{\code{code\_SPC\_ECoG\_STN\_Speech}}
(Vissani et al., 2025), run live on simulated coupled data so the exact method can be inspected.

\textbf{Synthetic data.} For each of $n_{\text{trials}}$ trials over $2$~s at $F_s = 500$~Hz
($N=1000$ samples), an LFP $x(t) = \sin(2\pi f t + \phi_0) + 0.25\,\xi$ is generated and its
instantaneous phase $\phi(t) = \arg\!\big(x(t) + \ii\,\mathcal H\{x\}(t)\big)$ is obtained from
the analytic signal via an FFT-based Hilbert transform. Spikes are an inhomogeneous Poisson
process: baseline probability $r_0/F_s$ per sample, multiplied inside a coupling window
$[\,\text{winC}\pm\text{winW}/2\,]$ by $\big(1 + \kappa\cos(\phi - \phi_{\text{target}})\big)$ with
\[
\kappa = \frac{\text{ES}-1}{\text{ES}+1}, \qquad
\text{ES} = \frac{1 + 2\sqrt{c}}{1 - 2\sqrt{c}},
\]
an effect-size parametrization of the coupling strength $c$.

\textbf{Coupling statistics.} For the $n$ spike phases in the window,
\[
\text{PLV} = \Big|\tfrac1n \textstyle\sum_k \E^{\ii\phi_k}\Big|, \qquad
\text{PPC} = \frac{n}{n-1}\Big(\text{PLV}^2 - \tfrac1n\Big),
\]
the pairwise phase consistency (Vinck et al., 2010) being an unbiased, rate-independent
estimator of coupling.

\textbf{Time-resolved analysis.} A constant-spike-count window (Fischer) de-biases the firing
rate: with target $=\max(40,\ \lfloor 0.1\,N_{\text{spk}}\rceil)$, each of $40$ time centres takes
the \emph{target} spikes nearest in time (first-crossing half-width). Significance comes from a
phase-shuffle permutation null --- each selected spike is re-paired with the phase drawn from a
random trial at the mirrored time index ($40$ permutations) --- and a center is significant when
its real PPC exceeds the $95^{\text{th}}$ percentile of the null. The preferred phase per center
is the circular mean of its selected spikes; it is overlaid on the across-trial
phase$\times$time density map only where the PPC is significant.

\textbf{Confidence interval.} The $95\%$ CI of the preferred (mean) phase follows the CircStat
routine \code{circ\_confmean} (Zar, eq.\ 26.24/26.25): with $r=\text{PLV}$, $R_n = n r$ and
$c_2 = \chi^2_{1,0.95}=3.841$, the half-width is $\delta = \arccos(t/R_n)$ where
\[
t =
\begin{cases}
\sqrt{\dfrac{2n\,(2R_n^2 - n c_2)}{4n - c_2}}, & r < 0.9,\\[10pt]
\sqrt{\,n^2 - (n^2 - R_n^2)\,\E^{c_2/n}}, & r \ge 0.9.
\end{cases}
\]
It is drawn as an arc just outside the polar plot.

\textbf{Controls.} \param{coupling $c$} $\in[0,0.23]$, preferred phase $\phi_{\text{target}}$,
oscillation frequency, base firing rate, number of trials, and the coupling-window
centre/width.

\vfill
\begin{center}\small\itshape
These notes are generated from the same source that runs the live demos; if a number here and
on the site ever disagree, the code is ground truth.
\end{center}

\end{document}
